\section*{第 \ref{ch:b1:continuity} 章——极限与连续}

\begin{solution}{exo:b1:continuity:1}
取 $u_n = \frac{1}{2\pi n + \pi/2}$ 与 $v_n = \frac{1}{2\pi n}$：
二者都趋于 $0^+$，但 $\sin\frac{1}{u_n} = 1$ 且
$\sin\frac{1}{v_n} = 0$。两个数列的像趋于不同的极限，
由\cref{thm:b1:continuity:seqchar}，在 $0^+$ 处无极限。

$x \sin\frac1x$：由 $\abs{x\sin\frac1x} \leq \abs x \to 0$ 夹逼，
故 $0$ 处极限存在且等于 $0$。
\end{solution}

\begin{solution}{exo:b1:continuity:2}
$f = \lfloor\cdot\rfloor$ 在 $\R \setminus \Z$ 上\omterm{def:b1:continuity:continuous}{连续}
（局部常值），在每个 $n \in \Z$ 处不\omterm{def:b1:continuity:continuous}{连续}：左极限
为 $n - 1$，函数值为 $n$。

$g(x) = x - \lfloor x\rfloor$：间断点与前述函数相同（恒等函数
\omterm{def:b1:continuity:continuous}{连续}，故 $g$ 继承 $f$ 的跳跃）；在 $n \in
\Z$ 处，左极限为 $1 \neq 0 = g(n)$。

$h(x) = \lfloor x\rfloor + (x - \lfloor x\rfloor)^2$：在 $\intco{n}{n
+ 1}$ 上 $h(x) = n + (x - n)^2$，此处\omterm{def:b1:continuity:continuous}{连续}；在 $x = n$ 处，左极限
为 $(n-1) + 1 = n = h(n)$，跳跃相互抵消。$h$ 在
$\R$ 上\omterm{def:b1:continuity:continuous}{连续}（且严格递增）。
\end{solution}

\begin{solution}{exo:b1:continuity:3}
$P(x) = x^5 - 3x + 1$：$P(-2) = -32 + 6 + 1 = -25 < 0$；$P(0) = 1 >
0$；$P(1) = -1 < 0$；$P(2) = 27 > 0$。三次变号分别出现在
互不相交的闭\omterm{prop:b1:reals:intervals}{区间} $\intcc{-2}{0}$、$\intcc{0}{1}$、$\intcc{1}{2}$ 上：
由\cref{thm:b1:continuity:ivt}，至少有三个根（
由于次数为 $5$，$P$ 至多有五个根；
通过单调性分析可知恰有三个。）
\end{solution}

\begin{solution}{exo:b1:continuity:4}
设 $P = a_{2m+1}X^{2m+1} + \dots$，其中 $a_{2m+1} > 0$（否则
将 $P$ 换为 $-P$）。提出主导项，$P(x) = a_{2m+1}
x^{2m+1}\bigl(1 + o(1)\bigr)$，$x \to \pm\infty$ 时，故 $P(x) \to
+\infty$（在 $+\infty$ 处），$-\infty$（在 $-\infty$ 处）。取 $a$ 使
$P(a) < 0$，取 $b$ 使 $P(b) > 0$，
由 $\intcc{a}{b}$ 上的介值定理知存在根。
\end{solution}

\begin{solution}{exo:b1:continuity:5}
设 $g(x) = f(x) - x$，它在 $\intcc{0}{1}$ 上\omterm{def:b1:continuity:continuous}{连续}。由于 $f$ 映入
$\intcc{0}{1}$，故 $g(0) = f(0) \geq 0$，$g(1) = f(1) - 1 \leq
0$。由\cref{thm:b1:continuity:ivt}，$g(c) = 0$ 对某个 $c$ 成立，即
不动点。

条件的必要性：在 $\intcc{0}{1}$ 上，间断\omterm{def:b1:logic:map}{映射}
$f(x) = 1$（当 $x \leq \frac12$ 时），$f(x) = 0$（当 $x > \frac12$ 时）无
不动点；在开\omterm{prop:b1:reals:intervals}{区间} $\intoo{0}{1}$（非闭\omterm{prop:b1:reals:intervals}{区间}）上，$f(x) =
\frac x2$ \omterm{def:b1:continuity:continuous}{连续}地映入 $\intoo{0}{1}$，且无不动点
（候选点 $0$ 不在其中）；在 $\R$ 上，$f(x) = x + 1$。
\end{solution}

\begin{solution}{exo:b1:continuity:6}
$g(x) = \cos x - x$ \omterm{def:b1:continuity:continuous}{连续}，$g(0) = 1 > 0$，$g(1) = \cos 1 -
1 < 0$，故方程在 $\intoo{0}{1}$ 内有解
（\cref{thm:b1:continuity:ivt}）。唯一性：$g$ 严格
在 $\R$ 上递减——当 $x \leq 0$ 时，$g(x) \geq 1 - x > 0$，本无零点；
且 $g'(x) = -\sin x - 1 \leq 0$，等号仅在
孤立点（$x \equiv -\frac\pi2 \bmod 2\pi$）处成立，故 $g$
严格递减（见\cref{ch:b1:derivative}；另一种证法：在
$\intcc{0}{1}$ 上，$\cos$ 严格递减，$-x$ 也是，故 $g$
亦然）。严格单调函数至多只有一个零点。
\end{solution}

\begin{solution}{exo:b1:continuity:7}
取定 $M = f(0) + 1$。存在 $A > 0$，使 $f(x) \geq M$，当 $\abs x
\geq A$ 时成立（由两个无穷极限的定义；取较大者为阈值）。
在闭\omterm{prop:b1:reals:intervals}{区间} $\intcc{-A}{A}$ 上，由极值定理
（\cref{thm:b1:continuity:evt}）给出 $c$，使 $f(c) =
\inf_{\intcc{-A}{A}} f \leq f(0)$。当 $\abs x \geq A$ 时，$f(x) \geq M
> f(0) \geq f(c)$。因此 $f(c)$ 为全局最小值。
\end{solution}

\begin{solution}{exo:b1:continuity:8}
在闭\omterm{prop:b1:reals:intervals}{区间} $\intcc{0}{T}$ 上，$f$ 有界，且能取到其上、下界
（\cref{thm:b1:continuity:evt}）；由周期性，这些界即为
$\R$ 上的界，且仍能取到。

设 $g(x) = f(x + \frac T2) - f(x)$，它\omterm{def:b1:continuity:continuous}{连续}。则
\[
g(0) + g\bigl(\tfrac T2\bigr)
= \bigl(f(\tfrac T2) - f(0)\bigr) + \bigl(f(T) - f(\tfrac T2)\bigr)
= f(T) - f(0) = 0 :
\]
$g(0)$ 与 $g(\frac T2)$ 异号（或其一为零），故
由 $\intcc{0}{T/2}$ 上的介值定理，存在 $c$ 使
$g(c) = 0$，即 $f(c + \frac T2) = f(c)$。
\end{solution}

\begin{solution}{exo:b1:continuity:9}
先证不等式：对 $0 \leq y \leq x$，
\[
\bigl(\sqrt y + \sqrt{x - y}\bigr)^2 = x + 2\sqrt{y(x-y)} \geq x,
\]
故 $\sqrt x \leq \sqrt y + \sqrt{x - y}$，即 $\sqrt x - \sqrt y
\leq \sqrt{x - y}$。于是 $\abs{\sqrt x - \sqrt y} \leq
\sqrt{\abs{x - y}}$ 对所有 $x, y \geq 0$ 成立。

\omterm{def:b1:continuity:uniform}{一致连续}性：给定 $\varepsilon > 0$，取 $\delta =
\varepsilon^2$；则 $\abs{x - y} \leq \delta$ 蕴含 $\abs{\sqrt x
- \sqrt y} \leq \sqrt\delta = \varepsilon$.

在 $0$ 附近不满足利普希茨条件：$\frac{\sqrt x - \sqrt 0}{x - 0} =
\frac{1}{\sqrt x} \to +\infty$（$x \to 0^+$），故不存在常数 $k$
能控制所有差商。
\end{solution}

\begin{solution}{exo:b1:continuity:10}
假设 $f$ 是\omterm{def:b1:logic:inj}{单射}、\omterm{def:b1:continuity:continuous}{连续}且不严格单调。
则存在 $a < b < c$ 且均属于 $I$，使 $f(b)$ 不介于 $f(a)$ 与
$f(c)$ 之间——事实上，若对\emph{所有}三元组，中间值都
介于两端值之间，则 $f$ 单调（任取两对比较；
简单的分情形检验即可）。不妨设 $f(b) > \max(f(a), f(c))$（另一情形
对称，将 $f$ 换成 $-f$）。取 $v$，使 $\max(f(a),
f(c)) < v < f(b)$。由介值定理，分别在
$\intcc{a}{b}$ 与 $\intcc{b}{c}$ 上，存在 $u_1 \in
\intoo{a}{b}$ 与 $u_2 \in \intoo{b}{c}$，使 $f(u_1) = v = f(u_2)$：
两个不同点的像相等，与\omterm{def:b1:logic:inj}{单射}性矛盾。
\end{solution}

\begin{solution}{exo:b1:continuity:11}
由 $f(0) = f(0+0) = 2f(0)$ 得 $f(0) = 0$；$f(-x) = -f(x)$ 由 $0 =
f(x - x)$ 得出。记 $\alpha = f(1)$。归纳得 $f(n) = n\alpha$，其中 $n
\in \N$，再由奇性对 $n \in \Z$ 成立。对 $q \in \N^*$：
$q\,f(\frac pq) = f(p) = p\alpha$（将 $\frac pq$ 自加 $q$
次），于是 $f(\frac pq) = \alpha\frac pq$，即 $f = \alpha\,
\mathrm{id}$ 在 $\Q$ 上成立。

设 $x \in \R$，$(r_n)$ 为有理数列，且 $r_n \to
x$（\omterm{def:b1:topology:dense}{稠密}性，见\cref{thm:b1:reals:density}，在嵌套
\omterm{prop:b1:reals:intervals}{区间}应用；或取 $r_n = \frac{\lfloor nx\rfloor}{n}$）。由\omterm{def:b1:continuity:continuous}{连续}性：
$f(x) = \lim f(r_n) = \lim \alpha r_n = \alpha x$.
\end{solution}

\begin{solution}{exo:b1:continuity:12}
设 $\varepsilon > 0$。由 $+\infty$ 处的极限，存在 $A$，使
$\abs{f(x) - \ell} \leq \frac\varepsilon2$ 对 $x \geq A$ 成立，于是对
$x, y \geq A$，有 $\abs{f(x) - f(y)} \leq \varepsilon$（
无需任何接近性条件）。

在线段 $\intcc{0}{A + 1}$ 上，由海涅定理
（见\cref{thm:b1:continuity:heine}）给出 $\delta_0 > 0$，对此
$\varepsilon$；令 $\delta = \min(\delta_0, 1)$。

任取 $x, y \geq 0$，满足 $\abs{x - y} \leq \delta$，不妨设 $x
\leq y$。若 $y \leq A + 1$，则两点都在该线段上，海涅
定理中的 $\delta_0$ 适用。否则 $y > A + 1$，于是 $x \geq y - 1 >
A$，两点都在 $\intco{A}{+\infty}$ 中，此时极限论证
适用。两种情形下均有 $\abs{f(x) - f(y)} \leq \varepsilon$，即
\omterm{def:b1:continuity:uniform}{一致连续}。
\end{solution}

\begin{solution}{pb:b1:continuity:1}
\textbf{1.} 对 $n \in \N$，由归纳法 $f(nx) = n f(x)$（
$f((n+1)x) = f(nx) + f(x)$）。又由 $f(0) = 2f(0)$ 得 $f(0) =
0$，且 $0 = f(x - x) = f(x) + f(-x)$ 给出奇性，于是 $f(nx) =
nf(x)$ 对 $n \in \Z$ 成立。对 $r = \frac pq$：$q\,f\bigl(\tfrac
pq x\bigr) = f(px) = p\,f(x)$，故 $f(rx) = r f(x)$，即 $f$ 是
$\Q$-线性的。

\textbf{2.} 由可加性，对所有 $x$ 与 $h$ 有：
\[
f(x + h) - f(x) = f(h) = f(x_0 + h) - f(x_0) .
\]
当 $h \to 0$ 时，右端趋于 $0$，由
$x_0$ 处的\omterm{def:b1:continuity:continuous}{连续}性；因此 $f(x + h) \to f(x)$，即在每点 $x$ 处\omterm{def:b1:continuity:continuous}{连续}。于是
由\cref{exo:b1:continuity:11}可得 $f(x) = f(1)\,x$。

\textbf{3.} 两个\omterm{def:b1:continuity:continuous}{连续}的可加函数具有
$cx$ 与 $c'x$ 的形式（第 2 小题）；若二者在\omterm{def:b1:topology:dense}{稠密}\omterm{def:b1:structures:subgroup}{子群}
$\Z + \sqrt2\,\Z$（\cref{exo:b1:reals:9}）上取值相同，则 $c\,g = c'g$
对该\omterm{def:b1:structures:subgroup}{子群}中某个 $g \neq 0$ 成立，于是 $c = c'$，两函数相等。无\omterm{def:b1:continuity:continuous}{连续}性时，
$\Q$-线性只把取值约束在
$\Q$-组合上，而 $1, \sqrt2$ 在 $\Q$ 上线性无关（$\sqrt2
\notin \Q$），因此 $f(1) = 0$、$f(\sqrt2) = 1$ 是相容的，并
在该\omterm{def:b1:structures:subgroup}{子群}上迫使
\[
f(m + n\sqrt2) = m\,f(1) + n\,f(\sqrt2) = n .
\]

\textbf{4.} 对 $t \in \intcc{0}{\ell}$，有 $a + t \in
\intcc{a}{b}$，故 $f(t) = f(a + t) - f(a) \leq M - f(a) =: M'$。

\textbf{5.} 对 $t \in \intcc{0}{\ell}$，同样 $\ell - t \in
\intcc{0}{\ell}$，由可加性，$f(t) = f(\ell) - f(\ell
- t) \geq f(\ell) - M'$。于是 $\abs{f} \leq C$ 在
$\intcc{0}{\ell}$ 上成立，其中 $C = \max\bigl(\abs{M'}, \abs{f(\ell) -
M'}\bigr)$。

\textbf{6.} 对 $t \in \intcc{0}{\ell/n}$，有 $nt \in
\intcc{0}{\ell}$，且 $f(t) = \frac{f(nt)}{n}$（第 1 小题），故
$\abs{f(t)} \leq \frac Cn$。给定 $\varepsilon > 0$，取 $n >
\frac C\varepsilon$：对 $0 \leq h \leq \frac\ell n$，
$\abs{f(h)} \leq \varepsilon$；对负的 $h$，利用 $f(h) =
-f(-h)$。于是 $f(h) \to 0 = f(0)$（当 $h \to 0$ 时），即在 $0$ 处\omterm{def:b1:continuity:continuous}{连续}，
从而处处\omterm{def:b1:continuity:continuous}{连续}（第 2 小题），于是 $f(x) = f(1)x$。

\textbf{7.} 若 $f$ 在 $\intcc{a}{b}$ 上单调不减，则
$f(a) \leq f(x) \leq f(b)$，故有界，由第 6 小题
可知 $f(x) = cx$；而 $c(b - a) = f(b) - f(a) \geq 0$ 迫使 $c
\geq 0$。

\textbf{8.} (a)$\Rightarrow$(b)$\Rightarrow$(c)：显然。
(c)$\Rightarrow$(a)：第 2 小题。(a)$\Rightarrow$(d)：线性
函数处处单调。(d)$\Rightarrow$(e)：单调
函数在 $\intcc{a}{b}$ 上被其
端点值所界。(e)$\Rightarrow$(a)：第 4--6 小题。五个
陈述等价——正则性阶梯。

\textbf{9.} 第 4 小题仅用到\omterm{def:b1:reals:bounds}{上界} $M$，第 5 小题
\emph{推出}了下界：由\omterm{def:b1:reals:bounds}{上界}借助
对称 $f(t) = f(\ell) - f(\ell - t)$。第 6 小题随后作用于
$\abs f \leq C$ 上。综上：可加且在一个非退化\omterm{prop:b1:reals:intervals}{区间}上\emph{从上方}有界
即蕴含线性。对于下界，
将此用于 $-f$（可加，且从上方有界）。最强
的一档 (e)：一个\omterm{prop:b1:reals:intervals}{区间}上的单侧界即已足够。

\textbf{10.} 若 $\frac{f(x)}{x}$ 恒为同一常数 $c$（对所有
$x \neq 0$），则 $f$ 为线性函数；于是存在非零的 $u, v$，
使 $\frac{f(u)}u \neq \frac{f(v)}v$，即 $u f(v) - v f(u)
\neq 0$：向量 $(u, f(u))$、$(v,
f(v))$ 的行列式非零，且二者张成 $\R^2$。

\textbf{11.} 对 $r, s \in \Q$，有 $f(ru + sv) = r f(u) + s f(v)$
（由第 1 小题两次及可加性），故图像含有
\[
\bigl(ru + sv,\; r f(u) + s f(v)\bigr)
= r\,(u, f(u)) + s\,(v, f(v)) .
\]
给定 $(x_0, y_0)$ 及 $\varepsilon > 0$：$2 \times 2$
方程组 $a(u, f(u)) + b(v, f(v)) = (x_0, y_0)$ 有唯一
实数解 $(a, b)$，因行列式非零。取
有理数 $r_n \to a$、$s_n \to b$，则 $r_n(u, f(u)) +
s_n(v, f(v)) \to (x_0, y_0)$ 按坐标收敛，而这些
点都在图像上：图像在 $\R^2$ 中\omterm{def:b1:topology:dense}{稠密}。

\textbf{12.} 设 $I$ 为非退化\omterm{prop:b1:reals:intervals}{区间}，$x_0$ 为其
中点，$M$ 任意：由\omterm{def:b1:topology:dense}{稠密}性，
以 $\min\bigl(\frac{\abs I}{2}, 1\bigr)$ 为界，图像上有点与 $(x_0, M + 1)$ 的距离不超过此界，即
$x \in I$ 且 $f(x) > M$，故在 $I$ 上无界，从而（由第 8 小题）
在每点都不\omterm{def:b1:continuity:continuous}{连续}，且在任何\omterm{prop:b1:reals:intervals}{区间}上都不单调；并且
对任意目标值 $y_0$，图像上靠近 $(x_0, y_0)$ 的点给出的 $f(x)$
可任意接近 $y_0$，其中 $x \in I$，即 $f(I)$ 在
$\R$ 中\omterm{def:b1:topology:dense}{稠密}。这是把第 8 小题倒过来读：既然各档
彼此等价，非线性的可加函数必然违反\emph{全部}
各档，且处处如此。

\textbf{13.} 良定义：由 $m + n\sqrt2 = m' + n'\sqrt2$ 推出
$(n - n')\sqrt2 = m' - m \in \Z$，故 $n = n'$（否则 $\sqrt2 \in
\Q$），从而 $m = m'$。$G$ 上的可加性
按坐标即显然。在 $0^+$ 附近无界：固定 $\varepsilon \in
\intoo{0}{1}$ 与 $N \in \N$。对每个固定的 $n$，满足 $\abs n
\leq N$ 时，条件 $m + n\sqrt2 \in \intoo{0}{1}$ 将 $m$ 限制在
长为 $1$ 的\omterm{prop:b1:reals:intervals}{区间}内，至多一个整数 $m$，对每个
$n$，故至多 $2N + 1$ 个属于 $G \cap \intoo{0}{1}$ 的元素满足
$\abs{f} \leq N$。但 $G \cap \intoo{0}{\varepsilon}$ 是
无限的（$G$ \omterm{def:b1:topology:dense}{稠密}，由 \cref{exo:b1:reals:9}），故其
中含有某个 $g$，使 $\abs{f(g)} > N$：$f$ 在
$0$ 的每个右\omterm{def:b1:topology:open}{邻域}内都无界。要把 $f$ 延拓为
$\R$ 上的可加函数，就是在一族实数上协调地取值，该族
实数 $\Q$-线性无关，且 $\R$ 在 $\Q$ 上由其张成
——即哈梅尔基；构造这样的基需要选择公理，
而任何显式公式都无法做到：构造性地，我们只能
在 $G$ 上掌握这一“怪物”。

\textbf{14.} $f(x) = f(\frac x2 + \frac x2) = f(\frac x2)^2 \geq
0$。若 $f(x_0) = 0$，则 $f(x) = f(x - x_0)f(x_0) = 0$ 对所有
$x$ 成立，此情形被排除。故 $f > 0$，且 $g = \ln \circ f$ \omterm{def:b1:continuity:continuous}{连续}
（由 \cref{prop:b1:continuity:algebra}），且 $g(x + y) = g(x) +
g(y)$：由 \cref{exo:b1:continuity:11}，$g(x) = cx$，于是 $f(x) =
\eu^{cx}$。反之，每个 $\eu^{cx}$ 都满足方程：\omterm{def:b1:continuity:continuous}{连续}
\omterm{def:b1:structures:morphism}{同态} $(\R, +) \to (\R^*, \times)$ 恰为
指数函数。

\textbf{15.} $g(u) = f(\eu^u)$ \omterm{def:b1:continuity:continuous}{连续}，且 $g(u + v) =
f(\eu^u \eu^v) = g(u) + g(v)$：$g(u) = cu$，且每个 $x > 0$
可写成 $x = \eu^u$，其中 $u = \ln x$：$f(x) = c\ln x$。

\textbf{16.} $h = \ln \circ f$ 在
$\intoo{0}{+\infty}$ 上\omterm{def:b1:continuity:continuous}{连续}，且 $h(xy) = h(x) + h(y)$：由第 15 小题，
$h(x) = c\ln x$，故 $f(x) = \eu^{c\ln x} = x^c$。

\textbf{17.} 取 $x = y = 0$：$f(0) = 2f(0) + f(0)^2$，故 $f(0)(1 +
f(0)) = 0$。若 $f(0) = -1$，令 $y = 0$，则 $f(x) = f(x) +
f(0) + f(x)f(0) = -1$ 对所有 $x$ 成立：常值函数 $f \equiv -1$
（它确实满足方程）。否则 $f(0) = 0$，记 $h
= 1 + f$ \omterm{def:b1:continuity:continuous}{连续}，$h(0) = 1$，并且
\[
h(x + y) = 1 + f(x) + f(y) + f(x)f(y) = h(x)\,h(y) :
\]
由第 14 小题，$h(x) = \eu^{cx}$，即 $f(x) = \eu^{cx} - 1$
（$c = 0$ 时 $f \equiv 0$）。全部解为：$f
\equiv -1$ 及 $f(x) = \eu^{cx} - 1$，$c \in \R$。

\textbf{18.} 取 $x = y = 0$：$2f(0) = 4f(0)$，故 $f(0) = 0$。取 $x =
0$：$f(y) + f(-y) = 2f(y)$，故 $f$ 为偶函数。取 $y = x$：$f(2x) =
4f(x)$。利用 $(x, y) \to (nx, x)$ 作归纳：
\[
f((n{+}1)x) = 2f(nx) + 2f(x) - f((n{-}1)x)
= (2n^2 + 2 - (n-1)^2) f(x) = (n+1)^2 f(x) .
\]
于是 $f(x) = f\bigl(q\cdot\frac xq\bigr) = q^2 f\bigl(\frac
xq\bigr)$ 给出 $f\bigl(\frac pq x\bigr) = \frac{p^2}{q^2}f(x)$：
$f(r) = r^2 f(1)$ 在 $\Q$ 上成立（由偶性处理符号）。两个
\omterm{def:b1:continuity:continuous}{连续}函数 $f$ 与 $x \mapsto f(1)x^2$ 在
\omterm{def:b1:topology:dense}{稠密}集 $\Q$ 上相等，故处处相等（由第 24 小题）：$f(x) =
f(1)\,x^2$。每个 $cx^2$ 都满足方程。

\textbf{19.} $g = f - f(0)$ \omterm{def:b1:continuity:continuous}{连续}，$g(0) = 0$，并且
满足延森方程（常数相消）。取 $y =
0$：$g\bigl(\frac x2\bigr) = \frac{g(x)}{2}$。于是对所有 $x,
y$:
\[
\frac{g(x + y)}{2} = g\Bigl(\frac{x+y}{2}\Bigr)
= \frac{g(x) + g(y)}{2} ,
\]
故 $g$ 可加且\omterm{def:b1:continuity:continuous}{连续}：$g(x) = cx$（由第 2 小题），并且
$f(x) = cx + d$，其中 $d = f(0)$。所有仿射函数都满足
延森方程，故所列即为全部解。

\textbf{20.} 对 $n$ 用归纳法。当 $n = 0$ 时：$\lambda \in \{0,
1\}$，显然。假设不等式对一切形如
$\frac{k}{2^n}$ 的权重成立。形如 $\lambda = \frac{k}{2^{n+1}}$ 的权重，其中 $k$
为偶数时归结为第 $n$ 层的情形；当 $k = 2j + 1$ 时，$\lambda$
为 $\lambda_1 = \frac{j}{2^n}$ 与 $\lambda_2 =
\frac{j+1}{2^n}$ 的中点。记 $z_i = \lambda_i x + (1 - \lambda_i)y$：
$\lambda x + (1-\lambda)y = \frac{z_1 + z_2}{2}$，故
\[
f\bigl(\lambda x + (1{-}\lambda)y\bigr)
\leq \frac{f(z_1) + f(z_2)}{2}
\leq \frac{(\lambda_1 + \lambda_2)f(x) + (2 - \lambda_1 -
\lambda_2)f(y)}{2}
= \lambda f(x) + (1 - \lambda)f(y) .
\]

\textbf{21.} 固定 $x, y$。\omterm{def:b1:logic:map}{映射} $\lambda \mapsto f(\lambda x +
(1 - \lambda)y)$ 与 $\lambda \mapsto \lambda f(x) + (1 -
\lambda)f(y)$ 在 $\intcc{0}{1}$ 上\omterm{def:b1:continuity:continuous}{连续}（复合与
代数运算，见\cref{prop:b1:continuity:algebra}）。不等式
在二进权重上成立，这些权重在 $\intcc{0}{1}$ 中\omterm{def:b1:topology:dense}{稠密}
（见\cref{exo:b1:reals:8}）；对任意 $\lambda$，取二进有理数
$\lambda_n \to \lambda$，并取极限
（用\cref{thm:b1:continuity:seqchar}与\cref{thm:b1:seq:order}）：
凸性不等式对每个 $\lambda \in
\intcc{0}{1}$ 成立——中点凸性加\omterm{def:b1:continuity:continuous}{连续}性即得凸性
（\cref{ch:b1:derivative}中的凸性概念）。

\textbf{22.} 非线性可加函数 $f$ 满足
$f\bigl(\frac{x+y}{2}\bigr) = \frac{f(x) + f(y)}{2}$ 恰好成立
（由第 1 题取 $r = \frac12$，再利用可加性）：它是
中点凸的，甚至是中点仿射的。若它满足完整的
凸性不等式（在某\omterm{prop:b1:reals:intervals}{区间} $\intcc{x}{y}$ 上），则对
$\lambda \in \intcc{0}{1}$ 有 $f(\lambda x + (1-\lambda)y) \leq
\max(f(x), f(y))$——于是在非退化\omterm{prop:b1:reals:intervals}{区间}上有\omterm{def:b1:reals:bounds}{上界}，
由第 9 题知其为线性，矛盾。可见第 21 题中的\omterm{def:b1:continuity:continuous}{连续}性
并非可有可无：没有它时，中点凸性
只控制可数的二进骨架，而
其间的\omterm{def:b1:continuity:continuous}{连续}统则会失控。

\textbf{23.} $p(x) = \frac{x^2}{2}$ 满足 $p(x+y) = p(x) +
p(y) + xy$。若 $f$ 为任一\omterm{def:b1:continuity:continuous}{连续}解，$g = f - p$ 是
\omterm{def:b1:continuity:continuous}{连续}且可加的，故 $g(x) = cx$：
\[
f(x) = \frac{x^2}{2} + cx , \qquad c \in \R ,
\]
且每个这样的函数都是解，因此列表是完整的。

\textbf{24.} 一般原理：若 $u, v$ \omterm{def:b1:continuity:continuous}{连续}且
在\omterm{def:b1:topology:dense}{稠密}集 $D \subseteq \R$ 上取值相同，则对 $x \in \R$ 取
$d_n \in D$，使 $d_n \to x$
（见\cref{prop:b1:topology:closureprops}）；$u(x) = \lim u(d_n) =
\lim v(d_n) = v(x)$。现设 $f$ 为
\omterm{def:b1:topology:dense}{稠密}\omterm{def:b1:structures:subgroup}{子群} $G$ 上的\omterm{def:b1:continuity:continuous}{连续}可加函数。固定 $x, y \in \R$，取 $g_n \to x$，
$g'_n \to y$，且 $g_n, g'_n \in G$；则 $g_n + g'_n \to x + y$
并由 $x + y$、$x$ 与 $y$ 处的序列\omterm{def:b1:continuity:continuous}{连续}性：
\[
f(x + y) = \lim f(g_n + g'_n) = \lim\bigl(f(g_n) +
f(g'_n)\bigr) = f(x) + f(y) :
\]
于是 $f$ 在整个 $\R$ 上可加（从而由第 2 题知其线性）。

\textbf{25.} (i) 对可加函数 $f$：线性 $\iff$ \omterm{def:b1:continuity:continuous}{连续}
$\iff$ 在某一点\omterm{def:b1:continuity:continuous}{连续} $\iff$ 在某\omterm{prop:b1:reals:intervals}{区间}上单调
$\iff$ 在某\omterm{prop:b1:reals:intervals}{区间}上有界（甚至单侧有界）。(ii) 图像在平面中
的\omterm{def:b1:topology:dense}{稠密}性表明，失效并非局部的
缺陷，而是全局的“爆炸”：在每个子\omterm{prop:b1:reals:intervals}{区间}之上，
函数值涂满整个 $\R$，因此每一种正则性质
处处同时失效。(iii) 要点：$cx$、$\eu^{cx}$、
$c\ln x$、$x^c$、$cx^2$ 与 $cx + d$——六种刻画，
同一方法：把方程化为柯西方程，先证明
$\Q$ 上的骨架（归纳法），再借\omterm{def:b1:topology:dense}{稠密}性
与\omterm{def:b1:continuity:continuous}{连续}性推广到 $\R$。(iv) $\Q$（或二进有理数）的\omterm{def:b1:topology:dense}{稠密}性在
\cref{exo:b1:continuity:11} 与第 21 题中承担了推广；它
在第 13 题中毫无作用，因为没有\omterm{def:b1:continuity:continuous}{连续}性时
值无法从\omterm{def:b1:topology:dense}{稠密}集传到其\omterm{def:b1:topology:closure}{闭包}——
\omterm{def:b1:topology:dense}{稠密}性只沿\omterm{def:b1:continuity:continuous}{连续}性传递信息。
\end{solution}
